Input to the Evapotranspiration (EVTA) Package is read from the file that has type ``EVT6'' in the Name File. Any number of EVT6 packages can be specified for a single groundwater flow model. All single-valued variables are free format.  The package type in the GWF Model name file PACKAGES block is always ``EVT6'' regardless of whether list-based or layer array-based input is used; the READASARRAYS option within the package input file activates layer array-based input.

Layer array-based input for evapotranspiration provides a similar approach for providing evapotranspiration rates as previous MODFLOW versions.  Layer array-based input for evapotranspiration can be used only with the DIS and DISV Packages.  Layer array-based input for evapotranspiration cannot be used with the DISU Package.

\mf will allow evapotranspiration in cells that also include a lake connection.  Refer to the ``\hyperref[sec:lakrchet]{Use of the Lake Package with RCH and EVT Packages}'' section for additional details when the EVTA and LAK packages are active in the same cell.

When layer array-based input is used for evapotranspiration, the DIMENSIONS block should not be specified.  The array size is determined from the model grid.   Segmented evapotranspiration cannot be used with the READASARRAYS option.

\vspace{5mm}
\subsubsection{Structure of Blocks}
\vspace{5mm}

\noindent \textit{FOR EACH SIMULATION}
\lstinputlisting[style=blockdefinition]{./mf6ivar/tex/gwf-evta-options.dat}
\vspace{5mm}
\noindent \textit{FOR ANY STRESS PERIOD}
\lstinputlisting[style=blockdefinition]{./mf6ivar/tex/gwf-evta-period.dat}
\packageperioddescriptionarray{evapotranspiration}

\vspace{5mm}
\subsubsection{Explanation of Variables}
\begin{description}
\input{./mf6ivar/tex/gwf-evta-desc.tex}
\end{description}

\vspace{5mm}
\subsubsection{Example Input File}
\lstinputlisting[style=inputfile]{./mf6ivar/examples/gwf-evta-example.dat}
